The increasing accessibility and falling cost of consumer drones has 
created a growing need for affordable detection systems that can be 
deployed by private individuals and small organisations. Existing 
commercial solutions based on radar and electro-optical sensors are 
priced far beyond the reach of most non-military users, while 
academic research in this area rarely reports hardware costs alongside 
detection performance. This thesis presents the design, implementation, 
and testing of a low-cost multimodal drone detection system built 
entirely from off-the-shelf and 3D-printed components, with a total 
hardware cost of approximately NOK~ 10,000.

The system combines a four-element MEMS microphone array with a 
pan-tilt camera platform running a YOLOv8n object detection model, 
following a coarse-to-fine detection strategy in which acoustic 
direction of arrival estimation orients the camera toward the target 
before visual detection takes over. Acoustic signal acquisition is 
handled by an STM32F413 microcontroller using the DFSDM peripheral 
for PDM-to-PCM conversion, with direction of arrival estimated using 
the GCC-PHAT algorithm. Sound classification is performed by two 
machine learning models: a Convolutional Neural Network and a 
decision-tree classifier, both trained on a combination of public 
datasets and data collected with the microphone array. The YOLOv8n 
model was trained on a public drone detection dataset and deployed 
on a Raspberry Pi~5 with a Hailo AI HAT+ accelerator, achieving 
real-time inference at approximately 108 frames per second.

The acoustic DOA estimation was demonstrated to accurately track a 
live drone circulating around the tower in a laboratory environment, 
with a full detection pipeline iteration time of approximately 1.285 
seconds. The CNN sound classifier achieved 99\% accuracy on 
microphone array recordings and produced no false positives or 
negatives during the live test, while the YOLOv8n model achieved an 
mAP@0.50 of 0.835 and an F1 score of 0.81 on the validation dataset. 
Real-time visual detection and MQTT coordinate publishing were 
successfully demonstrated in a live laboratory test. The full 
closed-loop camera tracking was not completed within the project 
timeframe due to issues with the motor control hardware.

The results demonstrate that a capable multimodal drone detection 
system can be realised at a fraction of the cost of commercial 
alternatives, and provide a foundation for future development toward 
a complete autonomous detection and tracking platform.


\newpage
\textbf{Norsk sammendrag}

Den økte tilgjengeligheten og de fallende prisene på forbrukerdrones 
har skapt et voksende behov for rimelige deteksjonssystemer som kan 
tas i bruk av privatpersoner og mindre organisasjoner. Eksisterende 
kommersielle løsninger basert på radar og elektro-optiske sensorer er 
priset langt over det de fleste ikke-militære brukere har råd til, 
mens akademisk forskning på dette feltet sjelden rapporterer 
maskinvarekostnader sammen med deteksjonsytelse. Denne oppgaven 
presenterer design, implementasjon og testing av et lavkost multimodalt 
dronedeteksjonssystem bygget utelukkende av hyllevare- og 
3D-printede komponenter, med en total maskinvarekostnad på 
omtrent NOK~10,000.

Systemet kombinerer et firekanalers MEMS-mikrofonarrray med en 
pan-tilt-kameraplatform som kjører en YOLOv8n 
objektdeteksjonsmodell, etter en grov-til-fin deteksjonsstrategi 
der akustisk ankomsretningsestimering orienterer kameraet mot 
målet før visuell deteksjon overtar. Akustisk signalanskaffelse 
håndteres av en STM32F413-mikrokontroller ved hjelp av 
DFSDM-periferiutstyret for PDM-til-PCM-konvertering, med 
ankomsretning estimert ved hjelp av GCC-PHAT-algoritmen. 
Lydklassifisering utføres av to maskinlæringsmodeller: et 
konvolusjonelt nevralt nettverk og en beslutningstreklassifikator, 
begge trent på en kombinasjon av offentlige datasett og data samlet 
inn med mikrofonarrayen. YOLOv8n-modellen ble trent på et offentlig 
dronedeteksjonsdatasett og distribuert på en Raspberry Pi~5 med en 
Hailo AI HAT+-akselerator, og oppnådde sanntidsinferens med 
omtrent 108 bilder per sekund.

Den akustiske DOA-estimeringen ble demonstrert å nøyaktig spore en 
levende drone som sirkulerte rundt tårnet i et laboratoriumsmiljø, 
med en fullstendig deteksjonspipeline-iterasjonstid på omtrent 
1,285 sekunder. CNN-lydklassifikatoren oppnådde 99\% nøyaktighet 
på mikrofonarrayopptak og produserte ingen falske positiver eller 
negativer under livetesten, mens YOLOv8n-modellen oppnådde en 
mAP@0.50 på 0,835 og en F1-score på 0,81 på valideringsdatasettet. 
Sanntids visuell deteksjon og MQTT-koordinatpublisering ble 
vellykket demonstrert i en levende laboratorietest. Den fullstendige 
lukkede sløyfe-kamerasporing ble ikke fullført innenfor 
prosjekttidsrammen på grunn av problemer med motorstyringsmaskinvaren.

Resultatene viser at et kapabelt multimodalt dronedeteksjonssystem 
kan realiseres til en brøkdel av kostnadene for kommersielle 
alternativer, og gir et grunnlag for videre utvikling mot en 
komplett autonom deteksjons- og sporingsplattform.







